\section{Schichtweise Entstehung mit dem Logit Lens}

Dieses Kapitel untersucht, ab welcher Schicht und wie deutlich sich positive und negative Eingaben innerhalb des Modells unterscheiden, ein zentraler Baustein für den beschreibenden Teil von Hypothese H3. Am Ende des Modells wird der finale Zustand mit der Ausgabematrix auf den Vokabularraum projiziert, so entstehen die Logits, aus denen das nächste Token vorhergesagt wird. Der \textbf{Logit Lens} wendet genau dieselbe Ausgabematrix bereits auf frühere, eigentlich noch unfertige Zwischenzustände an \cite{nostalgebraist2020logit}. So lässt sich für jede Schicht ablesen, welches Token das Modell zu diesem Zeitpunkt schon bevorzugen würde, ein Blick in den laufenden Denkprozess, wenn auch mit dem Vorbehalt, dass die Ausgabematrix nur für die letzte Schicht trainiert wurde und mittlere Schichten Information anders kodieren könnten.

Als Sentiment-Score dient die durchschnittliche Logit-Differenz zwischen allen positiven und allen negativen Hu-\&-Liu-Wörtern je Schicht. An einem Beispielpaar (\texttt{The food was tasty/disgusting and I feel very}) liegt dieser Wert in den frühen Schichten nahe null. Positive und negative Eingaben sind für das Modell an dieser Stelle kaum zu unterscheiden. Ab den mittleren Schichten baut sich eine Trennung auf, die zwischen Schicht 17 und 23 ihre höchsten Werte erreicht (Abbildung~\ref{fig:lddiff}). 
Bereits ab Schicht 4 zeigt sich eine erste Trennung zwischen positiven und negativen Eingaben. Der Abstand bleibt zunächst gering, wächst jedoch ab Schicht 17 deutlich an und erreicht in der letzten Schicht seinen stärksten Zuwachs. Dies deutet darauf hin, dass sentimentbezogene Informationen früh entstehen, aber erst in den späten Schichten klar ausgeprägt werden.

\begin{figure}[htbp]
  \centering
  \includegraphics[width=0.7\textwidth]{logit_lens_prompt_hu_liu_logit_difference_Trial2.png}
  \caption{Logit-Differenz je Schicht für ein Beispiel-Prompt-Paar, oben die Differenz je Prompt, unten die Lücke zwischen positivem und negativem Prompt.}
  \label{fig:lddiff}
\end{figure}



Damit dieses Muster nicht nur an einem einzelnen, handverlesenen Beispiel gilt, wird es über alle 1.707 Paare des CAD-Datensatzes aggregiert (Abbildung 5). Das bestätigt den Befund systematisch. Der mittlere Abstand liegt in den frühen Schichten zunächst nahe null. Ab etwa Schicht 4 ist jedoch eine erste leichte Abweichung von der Null-Linie erkennbar. Eine verlässliche Trennung besteht zu diesem Zeitpunkt noch nicht, da das Streuband weiterhin bis an die Null-Linie heranreicht und viele Prompt-Paare noch kaum unterscheidbar sind. In den darauffolgenden Schichten nimmt der Abstand zunächst nur langsam zu. Ab etwa Schicht 16 verstärkt sich dieser Trend jedoch deutlich, sodass die Trennung zwischen positiven und negativen Eingaben zunehmend ausgeprägter wird. Der größte Zuwachs zeigt sich schließlich im letzten Schritt von Schicht 23 auf 24, wo der mittlere Abstand sein Maximum erreicht.

\begin{figure}[htbp]
  \centering
  \includegraphics[width=0.85\textwidth]{logit_lens_cad_logit_difference_aggregate.png}
  \caption{Über 1.707 CAD-Paare gemittelte Logit-Differenz je Schicht (Mittelwert $\pm$ 1 Standardabweichung).}
  \label{fig:cad}
\end{figure}

Damit lässt sich festhalten, dass Sentiment nicht in einer einzelnen Schicht gespeichert, sondern schrittweise über den Residual Stream aufgebaut wird, mit der deutlichsten Trennung in den oberen Schichten 17 bis 23 und dem größten Einzelschritt am Übergang zur letzten Schicht. Wichtig ist dabei der methodische Vorbehalt aus der Einleitung dieses Kapitels. Da die verwendete Ausgabematrix nur für die letzte Schicht optimiert wurde, könnte der Logit Lens echte Sentiment-Information in mittleren Schichten unterschätzen, wenn diese dort noch in einer anderen, dem Modell eigenen Kodierung vorliegt und erst später in die vom Lens lesbare Form überführt wird. Dieser Befund macht den Bereich um Schicht 17 bis 23 trotzdem zum naheliegenden Kandidaten für die kausale Überprüfung in Kapitel 7; ob sich diese Erwartung bestätigt, zeigt sich dort.
